% ---------------------------------------------------------------------------
% exemplos_tema.tex — galeria de componentes do tema PMME/Insper
%
% NAO e uma apresentacao: e a referencia visual do tema, para quem for montar
% um deck novo (a banca 2, por exemplo) sem ter de garimpar exemplos dentro de
% banca1_beamer.tex. Cada frame mostra um componente com o codigo ao lado.
%
% Nao tem conteudo do projeto: todo texto aqui e placeholder. Por isso nao se
% sujeita as regras de proveniencia de figura e numero -- nao ha nenhum numero
% observado nem nenhuma figura de base do repositorio.
%
% Compilar SEMPRE a partir da raiz do repositorio:
%     bash scripts/apresentacao/build_exemplos_tema.sh
% ---------------------------------------------------------------------------
\documentclass[aspectratio=169,10pt]{beamer}
\usepackage{iftex}
\ifPDFTeX
  \usepackage[T1]{fontenc}\usepackage[utf8]{inputenc}\usepackage{lmodern}
\else
  \usepackage{fontspec}
\fi
\usepackage[brazil]{babel}
\usepackage{amsmath}
\usepackage{amssymb}
\usepackage{booktabs}
\usepackage{graphicx}
\usepackage{array}
\usepackage{tabularx}
\usepackage{ragged2e}
\usepackage{tikz}

\usetheme{PMME}

\title{Galeria de componentes}
\subtitle{Referência visual do tema PMME/Insper, para montar um deck novo}
\author{Projeto Mais Médicos Especialistas}
\orientacao{}
\institute{Insper}
\date{2026}
\pmmefontes{docs/07_apresentacoes/banca1/deck_beamer/fontes/}
\pmmeinsperlogo{docs/07_apresentacoes/banca1/deck_beamer/insper/logo.png}
\pmmeinsperfundo{docs/07_apresentacoes/banca1/deck_beamer/insper/insper-bg.png}
\pmmelogo{docs/07_apresentacoes/banca1/deck_beamer/insper/logo.png}
\pmmefoto{docs/07_apresentacoes/banca1/deck_beamer/assets/capa_foto.jpg}

\pmmeselecaoitem{Molduras}
\pmmeselecaoitem{Componentes}

\pmmesumariofundo{docs/07_apresentacoes/banca1/deck_beamer/assets/sumario_fundo.jpg}
\pmmesumariopasso{1.6}
\pmmesumarioitem{Molduras}{Capa, sumário, divisória e o rodapé com a trilha}
\pmmesumarioitem{Componentes}{O que entra no corpo de um frame}

% Fonte monoespacada miuda para as amostras de codigo desta galeria.
\newcommand{\cod}[1]{{\ttfamily\fontsize{6.4}{7.6}\selectfont\color{insperCinza}#1}}

\begin{document}

\begin{frame}[plain]\inspercapa\end{frame}
\begin{frame}[plain]\pmmesumario\end{frame}

% ===========================================================================
\section{Molduras}
\begin{frame}[plain]
  \pmmedivisoria{Molduras}{Capa, sumário, divisória e o rodapé com a trilha}
\end{frame}

\subsection{Moldura do frame}
\begin{frame}{A moldura: título, corpo, fontes e rodapé}
  \pmmebuild[o que a linha métrica faz aqui embaixo]{Título de build, com linha métrica}
  \vspace{-0.2em}
  {\footnotesize A moldura não tem barra de navegação no topo. A trilha de
  seções mora no rodapé, onde não custa altura de corpo: são \num{0,60 cm} de
  página, perto de uma linha e meia de texto a 10\,pt, que voltaram para o
  conteúdo. A seção corrente sai em negrito com o losango vermelho, seguida
  da subseção; as outras, em cinza.}

  \vspace{0.5em}
  {\footnotesize \lead{Registre a trilha no preâmbulo}, na ordem das seções,
  com os nomes curtos que o documento de conteúdo usa na linha de rastreio:}

  \vspace{0.35em}
  \cod{\textbackslash pmmeselecaoitem\{Motivação e Pergunta\}} \quad
  \cod{\textbackslash pmmeselecaoitem\{Literatura e Modelo\}} \quad
  \cod{\textbackslash pmmeselecaoitem\{Hipótese e Viabilidade\}}

  \vspace{0.5em}
  {\footnotesize Sem nenhum item registrado, o rodapé cai no
  \cod{\textbackslash rastreio\{...\}} manual, como antes.}

  \vspace{0.5em}
  {\footnotesize \lead{Depois de \cod{\textbackslash appendix}},
  \cod{\textbackslash pmmeapendice} troca a trilha pelo rótulo do apêndice e
  tira o contador, e o contador da parte principal para no total dela. A
  divisória aceita rótulo: \cod{\textbackslash
  pmmedivisoria[]\{Perguntas\}\{...\}} não põe número nenhum, \cod{[A]} põe a
  letra.}

  \fonte{\textbf{Fontes:} este é o bloco \cod{\textbackslash fonte\{...\}}, o
  único formato de fonte do deck.}
\end{frame}

% ===========================================================================
\section{Componentes}
\begin{frame}[plain]
  \pmmedivisoria{Componentes}{O que entra no corpo de um frame}
\end{frame}

\subsection{Equação em cartão}
\begin{frame}{Equação em cartão, com a referência na etiqueta}
  \pmmebuild{O componente}
  \vspace{-0.3em}
  \pmmeequacao[Autor et al. (2020), eq.~1, p.~184]{%
    \arg\max_{m \in M} \left\{ \sum_t \delta^t
    \left[ \frac{\mathbb{E}(w_{mt})}{p_{mt}} - c_{im} \right] \right\}}

  \vspace{-0.1em}
  \cod{\textbackslash pmmeequacao[Autor et al. (2020), eq.~1, p.~184]\{\textbackslash arg\textbackslash max ...\}}

  \vspace{0.7em}
  {\small Sem o argumento opcional sai só o quadro. A largura padrão é
  \num{0,94} da largura de texto e muda com \cod{\textbackslash
  pmmeequacaolargura\{...\}}; o respiro interno, com \cod{\textbackslash
  pmmeequacaofolga\{...\}}. A etiqueta põe a referência onde o olho já está,
  em vez de mandá-la para a linha de fontes.}
\end{frame}

\subsection{Fichas de termos}
\begin{frame}{Fichas de termos, no lugar da tabela de glossário}
  \pmmebuild{O componente}
  \pmmefichaaltura{2.4em}
  \begin{pmmefichas}[3]
    \pmmeficha{$m \in M$}{\textbf{municípios} candidatos}
    \pmmeficha{$\mathbb{E}(w_{mt})$}{remuneração esperada em $m$, no ano $t$}
    \pmmeficha{$p_{mt}$}{nível de preços local, o deflator}
    \pmmeficha{$c_{im}$}{custo \textbf{não pecuniário} de viver ali}
    \pmmeficha{$\sum_t \delta^t$}{a escolha é de \textbf{carreira}}
    \pmmeficha{$\bar{v}_i$}{a melhor alternativa do médico}
  \end{pmmefichas}
  \pmmefichaaltura{0pt}

  \vspace{0.5em}
  \cod{\textbackslash pmmefichaaltura\{2.4em\}}\\
  \cod{\textbackslash begin\{pmmefichas\}[3] \ \textbackslash pmmeficha\{\$m \textbackslash in M\$\}\{municípios candidatos\} ... \textbackslash end\{pmmefichas\}}

  \vspace{0.7em}
  {\small Duas ou três colunas. A altura fixa iguala os fundos quando as
  definições têm números de linhas diferentes; \num{0pt} volta à altura
  natural. Mesma informação que uma tabela de duas colunas, em menos altura.}
\end{frame}

\subsection{Selos e caixa de saída}
\begin{frame}{Selos de sinal e caixa de saída}
  \pmmebuild{Selos}
  {\small O selo dá o sinal sem obrigar a ler a expressão:
  \pmmeselov{eleva}\ \ \pmmeselor{reduz}\ \ \pmmeselo{obstáculo}}

  \vspace{0.4em}
  \cod{\textbackslash pmmeselov\{eleva\} \ \textbackslash pmmeselor\{reduz\} \ \textbackslash pmmeselo\{obstáculo\}}

  \vspace{1.0em}
  \pmmebuild{Caixa de saída}
  \pmmesaida{Gera $i^*$:}{a localidade que maximiza a renda real líquida
  esperada ao longo da carreira}

  \vspace{0.5em}
  \cod{\textbackslash pmmesaida\{Gera \$i\^{}*\$:\}\{a localidade que ...\}}

  \vspace{0.7em}
  {\small A caixa de saída nomeia o que a equação ou o bloco acima produz.
  Distingue-se de \cod{destaque}, que é a citação em destaque do documento de
  conteúdo, por ser alinhada à esquerda e trazer um rótulo.}
\end{frame}

\subsection{Layouts em partes}
\begin{frame}{Layouts em duas e em quatro partes}
  \begin{pmmeduas}
    \begin{pmmeparte}[a linha métrica também serve aqui]{Duas partes}
      \footnotesize \cod{pmmeduas} põe duas partes lado a lado, com o título
      em Playfair e o traço vermelho. As partes alinham pelo topo e a largura
      é calculada a partir da largura de texto.
    \end{pmmeparte}
    \begin{pmmeparte}{Quatro partes}
      \footnotesize \cod{pmmequatro} faz a grade 2$\times$2, com altura fixa
      de \num{0,36} da altura de texto e corpo em \cod{\textbackslash small}.
      Ambos aceitam \cod{[altura]}.
    \end{pmmeparte}
  \end{pmmeduas}

  \vspace{1.0em}
  \begin{pmmequatro}[0.2\textheight]
    \begin{pmmeparte}{Implementação}\footnotesize texto de exemplo\end{pmmeparte}
    \begin{pmmeparte}{Força de trabalho}\footnotesize texto de exemplo\end{pmmeparte}
    \begin{pmmeparte}{Capacidade}\footnotesize texto de exemplo\end{pmmeparte}
    \begin{pmmeparte}{Acesso}\footnotesize texto de exemplo\end{pmmeparte}
  \end{pmmequatro}
\end{frame}

\subsection{Véu com faixa}
\begin{frame}{Véu com faixa: a conclusão nasce da evidência}
  \pmmebuild{A evidência}
  {\small Este é o conteúdo do primeiro build. No segundo, ele continua na
  tela, desbotado, e a conclusão aparece numa faixa de sangria total por
  cima. É o recurso que o deck entregue na banca 1 usava para a pergunta e
  para a hipótese.}

  \vspace{0.6em}
  \begin{pmmeduas}
    \begin{pmmeparte}{Um lado}
      \footnotesize Texto de exemplo, para haver o que desbotar quando a
      faixa entrar no segundo build.
    \end{pmmeparte}
    \begin{pmmeparte}{O outro}
      \footnotesize Texto de exemplo, para haver o que desbotar quando a
      faixa entrar no segundo build.
    \end{pmmeparte}
  \end{pmmeduas}

  \vspace{0.7em}
  \cod{\textbackslash only<2>\{\textbackslash pmmeveuy\{6.5\}\textbackslash pmmeveu[0.93]\{A pergunta ...\}\textbackslash pmmeveubloco[5.25]\{...\}\}}

  \only<2>{%
    \pmmeveuy{6.2}%
    \pmmeveu[0.93]{A frase que fecha o argumento, em faixa de sangria total}%
    \pmmeveubloco[4.9]{%
      {\small O que mais entra no segundo build vai em \cod{\textbackslash
      pmmeveubloco}, porque o fluxo normal do frame continuaria abaixo do
      conteúdo do primeiro.}}%
  }
\end{frame}

\subsection{Imagem de fundo}
{\pmmefundo{docs/07_apresentacoes/banca1/deck_beamer/insper/insper-bg.png}
\begin{frame}{Imagem de fundo no frame}
  \pmmebuild{O componente}
  {\small A macro antes do \cod{\textbackslash begin\{frame\}} cobre a página
  inteira com a imagem, recortada e centrada. Vale até \cod{\textbackslash
  pmmesemfundo} ou até o fim do grupo, como aqui.}

  \vspace{0.6em}
  \cod{\{\textbackslash pmmefundo[0.18]\{caminho/da/imagem.png\}}\\
  \cod{\ \textbackslash begin\{frame\}\{Título\} ... \textbackslash end\{frame\}\}}

  \vspace{0.8em}
  {\small A opacidade padrão é \num{0,18}, que é a desta página: calibrada
  para o texto continuar legível por cima. Com opacidade \num{1} e frame
  \cod{[plain]}, é uma página de foto plena. A moldura fica sempre por cima
  do fundo.}
\end{frame}}

\end{document}
